\documentclass[10pt,a4paper]{amsart}
\usepackage[margin=19mm]{geometry}
\usepackage{amsmath,amssymb,mathtools}
\usepackage[scr=rsfs,cal=euler]{mathalfa}
\usepackage{xfrac}
\usepackage[unicode,hidelinks,bookmarks=false]{hyperref}
\newcommand{\ie}{i.e.\ }
\newcommand{\cf}{cf.\ }
\newcommand{\resp}{respectively\ }
\newcommand{\Subsection}{\subsection}
\providecommand{\nobreakdash}{\nobreak\hbox{-}}
\setlength{\emergencystretch}{2em}
\multlinegap0pt
\usepackage{amssymb}
\usepackage{xcolor}
\usepackage{booktabs}
\usepackage{multirow}
\usepackage{caption}
\usepackage{algorithm}\usepackage{algpseudocode}
\usepackage{natbib}
\usepackage[shortlabels]{enumitem}
\usepackage{mathtools}
\usepackage{siunitx}
\usepackage{tikz}
\newtheorem{proposition}{Proposition}
\newcommand{\capacity}{b}
\newcommand{\firstDP}{{\tt DP_1}}
\newcommand{\profit}{p}
\newcommand{\weight}{w}
\title{The colored knapsack problem: structural properties and exact algorithms}
\author{Fabio Ciccarelli, Alexander Helber, Erik M\"uhmer}
\date{Source: arXiv:2602.12214v1 (CC BY 4.0)}
\begin{document}
\maketitle

\noindent\emph{Source and licence.} The colored knapsack problem: structural properties and exact algorithms. Version arXiv:2602.12214v1 (CC BY 4.0), distributed by arXiv under the Creative Commons Attribution 4.0 International licence, http://creativecommons.org/licenses/by/4.0/, as recorded in the arXiv licence metadata for this exact version and shown on the abstract page https://arxiv.org/abs/2602.12214v1. Full licence notice: \texttt{licenses/arxiv-2602.12214-CC-BY-4.0.txt}.\par\smallskip

\noindent\textit{Editorial scope.} Counted unit: the article's proposition proposition (source0.tex lines 578--580). The article's complete original proof of it is delivered with it (source0.tex lines 581--629, one \texttt{proof} environment of 49 lines). No mathematics is written, repaired or re-derived: the statement and the proof are the article's own lines, in the article's own notation, and no retained line is substituted or re-flowed.  Retained uncounted as the source-local context the proof needs: lines 531--533.  Nothing was omitted from any retained span, so the package carries no omission record.  No retained line contains a citation, so the wrapper carries no bibliography and no bibliography entry.  No retained line contains a figure environment, an image-inclusion command or drawing code, so no image is delivered and the package declares no asset dependency.  Editorial formatting only, in addition: an AMS \texttt{amsart} preamble that restates the article's own declarations and macros used by the retained lines, namely the environments \texttt{proposition} on separate counters, together with the standard packages those lines need.  The retained environments print in this document as Proposition 1 in order of appearance, whereas the article prints the same items with the numbering of its own full document.  The complete native source of the article as distributed in the arXiv e-print for this exact version is bundled unchanged as \texttt{sources/194531/source0.tex}.
\begin{equation}
    f(i,t,d,a,q) = \max\bigg\{\profit(S) \mid S \subseteq \{1,2,\dots, i\}, \, |S| = t, \, \overline{k}(S) = d, \, k_{\kappa_i}(S) = a, \, \weight(S) = q\bigg\}, \label{dp1:goal}
\end{equation}

\begin{proposition}
The $\firstDP$ algorithm computes the value of $f(i,t,d,a,q)$ for all $i,t,d,a \in \{0,1,\dots,n\}$, $q \in \{0,1,\dots, \capacity\}$ according to~\eqref{dp1:goal}, if a corresponding solution exists, in at most $\mathcal{O}(\capacity \, n^4)$ operations.
\end{proposition}

\begin{proof}
For $i \in \{0,1,\dots,n\}$, $t,d,a \in \{0,1,\dots,n\}$, and $q \in \{0,1,\dots,\capacity\}$, let $f^*(i,t,d,a,q)$ denote the maximum profit obtainable by selecting a subset of the first $i$ items such that exactly $t$ items are selected, the maximum number of selected items of any color equals $d$, exactly $a$ selected items have color $\kappa_i$, and the total weight equals $q$.
If no such subset exists, then we assume $f^*(i,t,d,a,q) = -\infty$.
The $\firstDP$ algorithm maintains a table $f(i,t,d,a,q)$, and we prove that $f(i,t,d,a,q) = f^*(i,t,d,a,q)$ for all $i,t,d,a,q$ by induction on $i$.

\smallskip

For $i=0$, no items are available. The only feasible solution is the empty set, which satisfies $t=d=a=q=0$, and has profit $0$. 
Thus, $f^*(0,0,0,0,0) = 0$, and all other states are infeasible. The algorithm correctly initializes this state and assigns $-\infty$ to all others, so the claim holds for $i=0$.

\smallskip

We now assume that, for some $i \ge 1$, the equality $f(i-1,t,d,a,q) = f^*(i-1,t,d,a,q)$ holds for all $(t,d,a,q)$. We therefore want to show that $\firstDP$ correctly computes $f(i,\cdot)$.
The algorithm iterates over all reachable states $(t,d,a,q)$ at stage $i-1$, i.e., those with $f(i-1,t,d,a,q) > -\infty$, and determines the correct counter $a'$ for the current color $\kappa_i$. 
If $i=1$ or $\kappa_i = \kappa_{i-1}$, then the counter is preserved, so $a' = a$. 
Otherwise, since items are sorted by color, item $i$
is the first occurrence of color $\kappa_i$, and the counter is correctly reset to $a' = 0$. 
From each such state, it considers the two exhaustive and mutually exclusive choices: excluding or including item $i$ among the selected items. 
It follows that: 
i) if $i$ is not selected, the selected set remains unchanged. The total number of selected items, the number of items of each dominant color, and the total weight remain $t$, $d$, and $q$, respectively. Only the counter for the current color is updated to $a'$. The resulting state is therefore $(i,t,d,a',q)$. The algorithm updates $f(i,t,d,a',q)$ by taking the maximum between its current value and the propagated value $f(i-1,t,d,a,q)$; 
ii) if item $i$ is selected, the total number of selected items increases to $t+1$, the total weight to $q+\weight_i$, and the counter of the current color to $a'+1$. 
The value of $d$ must be updated accordingly. 
Since $d$ is the maximum frequency among all colors in the partial solution at stage $i-1$, and the current color count was $a'$, the new dominant frequency is $\max\{d,\, a'+1\}$. 
Because $d \ge a'$, this update can be written as $d + \sigma(a',d)$, where
$\sigma(a',d) = 1$ if $a' = d$ and $0$ otherwise. 
The algorithm thus transitions to the state $(i,\, t+1,\, d+\sigma(a',d),\, a'+1,\, q+\weight_i)$ and updates its value to $f(i-1,t,d,a,q) + \profit_i$, taking the maximum with any previously stored value.

\smallskip

To show that these transitions preserve optimality, consider an optimal solution to a given state $(i,t,d,a,q)$, of value $f^*(i,t,d,a,q)$. 
Item $i$ is either selected or not. 
If item $i$ is not selected, the solution uses only items from $\{1, 2,\dots,i-1\}$ and corresponds to a feasible state at stage $i-1$ with the same profit. 
By the inductive hypothesis, this value is stored in $f(i-1,\cdot)$ and is propagated to stage $i$ by the ``do not pack'' transition.
If item $i$ is instead selected, removing it yields a feasible solution for the first $i-1$ items. 
This partial solution must be optimal for its parameters; otherwise, replacing it with a better one would improve the full solution, contradicting optimality. 
By the inductive hypothesis, its value is correctly stored in $f(i-1,\cdot)$, and the ``pack'' transition adds
$\profit_i$ and updates the state parameters consistently.\\
Since the algorithm considers both cases for every reachable state at stage $i-1$ and updates each destination state by taking the maximum profit over all incoming transitions, it follows that $f(i,t,d,a,q) = f^*(i,t,d,a,q)$ for all $t,d,a,q$.
By induction, this holds for all $i \in \{0,1,\dots,n\}$.

\smallskip

The state space of $\firstDP$ is defined by
$i,t,d,a \in \{0,1,\dots,n\}$ and $q \in \{0,1,\dots,\capacity\}$, and thus contains
$\mathcal{O}(n^4 \capacity)$ states. 
For each reachable state, the algorithm performs a
constant number of operations. 
Consequently, the overall time complexity of $\firstDP$ is $\mathcal{O}(n^4 \capacity)$.
\end{proof}

\end{document}
